\section{Discussion}
\paragraph*{What the evidence supports.} In simulation on the unvalidated twin, a residual PPO around a validation-tuned constant beat that constant by about $0.5$--$0.6$ median return, consistently across three training seeds and on fresh scenarios. But the same data show that most of this gain is available without learning. A larger validation search for the constant recovers $\SensVThreeStaticBestConstantExtBestConstantMed$, $\SensVThreeChangesBestConstantExtBestConstantMed$ and $\SensVThreeOodBestConstantExtBestConstantMed$ of it in the three regimes, and a non-learned rule with three load features matches the policy within noise. The policy's gain over the stronger constant is small ($\SensVThreeStaticPpoBestConstantExtMed$ static, $\SensVThreeChangesPpoBestConstantExtMed$ with changes, $\SensVThreeOodPpoBestConstantExtMed$ out of distribution) and in the payload-change regime not distinguishable from zero. The supported conclusion is therefore protocol-level: \emph{an apparent RL gain over a modestly tuned constant can mostly be a better-tuned constant}, and papers should report a validation-tuned constant with a generous search and a non-learned scenario-conditioned baseline next to any learned tuner.

\paragraph*{Does the policy adapt at all?} Its gains are not constant: the scatter of final gains against scenario parameters (Fig.~\ref{fig:scatter}) shows that $K_p$ is markedly higher when gravity acts (about $27$--$37$ against $22$--$27$), i.e.\ the policy tracks the gravity load, in the way the non-learned rule does. It learned a plausible, simple dependence on the load; it did not learn more than the simple rule captures.

\paragraph*{Why the unconstrained setup failed (hypotheses, not shown).} v2 differed from v3 in four ways at once (unconstrained actions, an under-shaped reward, a checkpoint chosen on the training-evaluation stream, early-peaking training), and the experiments do not isolate which mattered. The gains piled up at the bounds, suggesting that without a reference the policy exploits the cost structure at the edges of the action range.

\paragraph*{Mean versus median.} Medians and means of the paired differences can disagree: on the confirmatory seeds the one-shot bandit is slightly below the best constant in the median in the static regime ($\ConfThreeStaticBanditMed$) but above it in the mean ($\ConfThreeStaticBanditMean$), because a minority of scenarios with large gains dominate the mean. This is why median, mean and interquartile range are reported side by side and why a median interval alone should not be read as an effect size for the mean.

\paragraph*{Headroom.} The hindsight oracle is far ahead of every deployable method (the policy captured $\ConfThreeStaticCapMean$, $\ConfThreeChangesCapMean$ and $\ConfThreeOodCapMean\,\%$ of its headroom on average over seeds), so there is room for scenario-aware tuning; the evidence here is that neither this PPO configuration nor a simple load-conditioned rule reaches it.

\paragraph*{Practical guidance.} Report (i) a validation-tuned constant with a search budget at least as large as the learner's, (ii) a non-learned scenario-conditioned baseline, (iii) a hindsight oracle, (iv) infeasible scenarios separately, and (v) paired differences with median, mean and spread. Treat success-rate gains with suspicion unless they replicate across seed sets.
